\section{Related Work}
\label{sec:related}

Agent memory systems can be organized along three axes: (i) what is stored (raw turns, extracted facts, structured knowledge), (ii) how storage is organized (vector index, hierarchical tiers, knowledge graph), and (iii) the deployment surface (library, OS service, network protocol). \Cref{tab:related} situates AMP and prior work along these axes.

\paragraph{Long-Context Language Models.} The simplest baseline for long-horizon recall is to fit history directly in the model's context window. This approach scales poorly: cost grows with history, and attention degrades non-uniformly across position. \citet{liu2024lost} document a ``lost in the middle'' effect in which information placed in the middle of long contexts is recovered substantially worse than information at the ends, even for models explicitly designed for long contexts. Memory-as-context, in other words, is not free of representational pathologies.

\paragraph{Retrieval-Augmented Generation.} \citet{lewis2020rag} introduced RAG, in which a non-parametric vector index is queried at generation time to inject relevant passages into the prompt. RAG is the dominant pattern for static knowledge but treats every chunk as independent. Applied naively to multi-session dialogue, the chunking step destroys turn-level structure: speaker identity, temporal ordering, and the dependency of each utterance on its predecessors are flattened into bag-of-text. We include naive RAG as a baseline in Section~\ref{sec:experiments} for exactly this reason.

\paragraph{Extraction-First Memory Systems.} Mem0~\citep{mem0} extracts and consolidates ``salient information'' from ongoing conversations into a structured memory store, optionally enhanced with a knowledge graph. The published results report a 26\% LLM-judge improvement over a full-context baseline and >90\% token cost savings on the LoCoMo benchmark. MemoryBank~\citep{memorybank} uses a similar extract-and-retrieve pattern with an Ebbinghaus-curve-inspired forgetting mechanism applied to memory items. We refer to this family as \emph{extraction-first}: the conversational record is summarized into atoms before being stored. The pattern is efficient, but the compression is irreversible --- once an utterance is paraphrased into a fact, the surrounding turn-level context is gone. Section~\ref{sec:experiments} characterizes when this matters.

\paragraph{Operating-System-Style Memory.} MemGPT~\citep{memgpt} treats LLM context as a memory hierarchy in the operating-system sense: a small ``main'' context for the model, larger archival storage paged in and out under explicit control. AMP shares the tiered framing but differs in three respects: (i) AMP's STM preserves verbatim turns rather than paged context summaries, (ii) AMP adds an explicit consolidation step that yields semantically indexed insights rather than file-system-style archival, and (iii) AMP's interface is a network protocol (MCP) rather than an in-process Python framework. Letta is a continuation of MemGPT as a deployable service; we cite the foundational paper.

\paragraph{Graph-Based Agent Memory.} Zep~\citep{zep} introduces Graphiti, a temporally-aware knowledge graph for agent memory, and reports state-of-the-art results on Deep Memory Retrieval (94.8\%, vs MemGPT 93.4\%) and substantial gains on LongMemEval. A-MEM~\citep{amem} draws on the Zettelkasten note-linking method to evolve a network of structured memory notes via dynamic indexing. Both systems treat the graph as the primary memory substrate. AMP's entity layer is intentionally minimal: it stores extracted entities as \emph{nodes} only (no edges, no co-occurrence weights), on the hypothesis that turn-level recall through STM and LTM is the load-bearing component for conversational benchmarks. A richer graph layer with weighted edges is left as future work; we do not claim a graph-based contribution.

\paragraph{Side-Network Memory.} LongMem~\citep{longmem} couples a frozen LLM backbone to an adaptive side network that retrieves from a 65k-token external memory cache. This trades architectural changes for context-window extension. AMP is purely a system-level approach: it makes no modifications to the underlying LLM, which is essential for compatibility with closed-weight models accessed over a protocol.

\paragraph{Positioning.} AMP combines (i) verbatim short-term retention, (ii) consolidation-based long-term insights with recency-weighted vector retrieval, and (iii) a Model Context Protocol surface, under a local-first deployment model. The combination of MCP-native deployment with cognitive-architecture-inspired tiering is, to our knowledge, the first such proposal in the agent memory literature. Section~\ref{sec:experiments} provides empirical evidence that this combination is competitive with naive RAG, narrowly trails Full Context, and significantly outperforms an extraction-first baseline on the LoCoMo benchmark at our evaluation scale.
